\chapter{Giới thiệu}
\label{ch:gioithieu}

\section{Bối cảnh và động cơ}

Phát hiện biển báo và đèn tín hiệu giao thông là bài toán nền tảng của hệ hỗ trợ lái và xe tự
hành. Với các kiến trúc hiện đại như YOLOv8, độ chính xác trên các bộ dữ liệu chuẩn đã rất cao:
notebook nguồn của đề tài đạt mAP@0,5 = 0,970 trên tập kiểm thử. Khi con số tổng đã bão hoà, câu
hỏi có giá trị khoa học không còn là ``mô hình nào tốt hơn'' mà là ``con số đó có \emph{đáng tin}
và \emph{đồng đều} không''.

\section{Định vị đề tài --- hai hướng tiếp cận}

Có hai hướng tiếp cận khi làm việc với một bài toán học máy đã có kết quả tốt. Hướng
\emph{model-centric} tập trung cải tiến mô hình để đẩy chỉ số lên cao hơn. Hướng
\emph{data-centric} \cite{datacentric} lại xem chất lượng và cấu trúc dữ liệu là đối tượng
nghiên cứu, và đặt câu hỏi về \emph{độ tin cậy} của chính con số đã đạt được. Đề tài này đi theo
hướng data-centric, như trình bày ở Bảng~\ref{tab:phanbiet}.

\begin{table}[H]
\centering
\caption{Hai hướng tiếp cận với một bộ phát hiện đã huấn luyện}
\label{tab:phanbiet}
\small
\begin{tabularx}{\linewidth}{@{}p{3.0cm}XX@{}}
\toprule
 & \textbf{Hướng model-centric} & \textbf{Hướng data-centric (đề tài này)} \\
\midrule
Câu hỏi & Mô hình nào phát hiện tốt hơn? & mAP 0,97 có đáng tin và đồng đều không? \\
Vai trò mô hình & Đối tượng nghiên cứu --- huấn luyện, so sánh & Công cụ đo cố định --- chỉ phân
tích lại kết quả \\
Đầu ra & Bảng benchmark, mAP tổng & Bằng chứng thống kê về độ tin cậy và phân hoá lớp \\
Phương pháp & Thiết kế, huấn luyện mạng & Kiểm định thống kê, phân tích rò rỉ và phân hoá \\
\bottomrule
\end{tabularx}
\end{table}

Mô hình YOLOv8s đã huấn luyện được xem là công cụ đo cố định; đề tài chạy ở chế độ
\textit{results-only} --- không huấn luyện lại mà phân tích lại các kết quả nó sinh ra. Trọng số
YOLOv8 đã huấn luyện (\texttt{best.pt}) được đính kèm trong thư mục \texttt{models/} của đề tài
để phục vụ ứng dụng demo và tái lập.

\section{Câu hỏi nghiên cứu}

\begin{itemize}[leftmargin=1.4cm,label=]
  \item \textbf{CH-A.} Mức độ rò rỉ dữ liệu xuyên split là bao nhiêu, và nó ảnh hưởng thế nào tới
  độ tin cậy của mAP?
  \item \textbf{CH-B.} Hiệu năng có đồng đều giữa các lớp không? Nếu không, khoảng cách lớn nhất
  nằm ở đâu?
  \item \textbf{CH-C.} Khoảng cách đó do độ hiếm của lớp (mất cân bằng) hay do bản chất vật thể?
\end{itemize}

\section{Đóng góp}

\begin{itemize}
  \item Định lượng phơi nhiễm rò rỉ dữ liệu và chỉ ra khoảng trống trong quy trình khử trùng của
  notebook nguồn.
  \item Chứng minh bằng kiểm định thống kê rằng đèn tín hiệu kém hơn biển báo một cách có ý nghĩa.
  \item Bác bỏ giả thuyết ``mất cân bằng gây ra thất bại'' bằng bằng chứng tương quan, chuyển
  nguyên nhân sang loại vật thể.
  \item Cung cấp pipeline tái lập, ứng dụng web, báo cáo Word/LaTeX và slide trình bày.
\end{itemize}
